\chapter{From Source Files to Submission}
\label{chap:submission}
\index{submission}\index{metadata}\index{font embedding}\index{accessibility}\index{manifest}

Finishing is a controlled transition from an editable research project to an
exact release. The final PDF, source archive, supplementary files, and portal
metadata must describe the same work. No universal checklist replaces the
current instructions of a journal, institution, funder, repository, or client.

\begin{learningoutcomes}
  \item Freeze a candidate release and apply a venue-specific compliance matrix.
  \item Proofread source, PDF, references, citations, equations, tables, and figures.
  \item Inspect metadata, fonts, page properties, links, and accessibility needs.
  \item Prepare clean source, supplementary files, a manifest, and a final PDF.
  \item Archive an editable reproducible project and the exact submitted release.
  \item Conduct a final submission rehearsal and record evidence.
\end{learningoutcomes}

\section{Define the release candidate}

Choose a version that has completed scientific review, author approval, and the
required ethics or organisational checks. Tag it or record an immutable commit
identifier. During final quality control, accept only deliberate corrections;
unbounded rewriting restarts review.

Create a release register with:

\begin{itemize}
  \item document title, version, date, and main-file path;
  \item venue and manuscript or deposit type;
  \item source-version identifier;
  \item responsible author and approver;
  \item required files and submission deadline;
  \item status of scientific, editorial, technical, and permission checks.
\end{itemize}

\begin{submissioncheck}
Download or open the current official instructions at the start of the finishing
stage. Record their version or access date. Recheck them immediately before
upload because file, declaration, and metadata requirements can change.
\end{submissioncheck}

\section{Build a compliance matrix}

Turn every instruction into a check with evidence:

\begin{center}
\begin{tabularx}{\textwidth}{@{}p{.21\textwidth}Xp{.22\textwidth}@{}}
  \toprule
  Requirement & Implementation & Verification evidence \\
  \midrule
  Anonymous main file & Anonymous build variant & Identity audit \\
  Figure formats & Separate accepted files & File inspection \\
  Reference style & Venue class or style & Sample comparison \\
  Data statement & Declarations section and form & Author approval \\
  Page limit & Final generated PDF & PDF page count \\
  \bottomrule
\end{tabularx}
\end{center}

Write ``not applicable'' only with a reason. An unresolved instruction should
be escalated to the appropriate editorial office, graduate office, repository,
or client rather than settled by appearance.

\section{Proofread at three levels}

\subsection{Scientific and logical review}

Check the question, methods, analyses, results, limitations, and conclusions as
a chain. Confirm sample sizes, denominators, exclusions, equation definitions,
units, uncertainty, statistical directions, and the difference between
association and causation. Trace each headline claim to evidence.

\subsection{Source review}

Search the source for draft markers, manual numbering, broken labels, duplicate
metadata, obsolete files, comments not intended for release, and private notes.
Inspect changes since the last approved version. Source review can find a hidden
conditional or unused old section that the PDF alone cannot reveal.

\subsection{Rendered-PDF review}

Read the PDF in page order, ideally on screen and in a print-like view. Check
page breaks, widows and orphans where severe, clipped content, running heads,
blank pages, table continuation, float order, caption proximity, footnotes,
bookmarks, and link destinations. Compare the table of contents and lists with
the body.

\begin{pausepredict}
Why can source spell-checking not detect a figure label clipped outside the
page? Why can PDF reading not reveal an author name hidden in an unused source
file included in the submission archive? Assign each issue to the appropriate
review layer.
\end{pausepredict}

\section{Cross-references and citations}

Search the final log for undefined references and citations. Click a sample of
internal and external links, including entries near page boundaries. Confirm
that each ``above'', ``below'', table number, figure number, equation number,
chapter, and appendix still points correctly after pagination changed.

For citations:

\begin{itemize}
  \item verify that each source supports the attached claim;
  \item compare author, title, year, venue, pages, DOI, and edition with an
        authoritative record;
  \item distinguish preprints, accepted manuscripts, versions of record, data,
        software, standards, and personal communications;
  \item remove uncited entries if the venue expects only cited works;
  \item check that retracted, corrected, or superseded sources are handled
        appropriately.
\end{itemize}

Never fabricate a DOI, page range, or access date to make a record look complete.
Use a visible placeholder until the fact is verified.

\section{Equations and scientific notation}

Read equations independently from the prose. Check:

\begin{itemize}
  \item every symbol is defined before or immediately after use;
  \item vectors, matrices, variables, operators, and units use consistent style;
  \item signs, indices, parentheses, limits, and equation directions match the
        analysis;
  \item displayed equations fit at final column width;
  \item numbering and punctuation support the surrounding sentence;
  \item significant figures and rounding are consistent with the evidence.
\end{itemize}

Have a domain expert verify equations whose visual plausibility can conceal a
scientific error. \LaTeX{} cannot know that a denominator is wrong.

\section{Tables}

Audit every table against its generating output or source record. Check row and
column order, headings, units, denominators, missing-value notation, rounding,
footnotes, statistical definitions, and totals. Confirm that the caption and
notes make the table intelligible without searching the prose.

Inspect tables at the final page and column width. Rules should not cross text;
decimal alignment should not hide signs; symbols must be defined; repeated
header rows should appear on multi-page tables. A screenshot of a spreadsheet
is not an editable, accessible substitute for a real table.

\section{Figures and image preflight}

For each figure, verify content and file properties:

\begin{enumerate}
  \item compare plotted values with the final analysis;
  \item confirm axes, units, legends, uncertainty, sample size, and panel labels;
  \item state whether data are simulated, illustrative, or empirical;
  \item inspect colour contrast and grayscale legibility;
  \item check text at intended physical size, not only when zoomed;
  \item preserve vector graphics for lines and text where accepted;
  \item verify raster pixel dimensions and the venue's effective-resolution rule;
  \item record creator, source, licence, permission, and required credit.
\end{enumerate}

Effective resolution depends on placed size. A 1200-pixel image placed at
\SI{4}{\inch} provides 300 pixels per inch; enlarging it to
\SI{8}{\inch} provides only 150. Do not claim universal resolution thresholds;
follow the venue's current specifications for each image type.

\begin{scienceinpractice}
A microscopy plate may require an unaltered original archive, a documented
processing history, scale bars derived from calibration, panel-level labels,
and separate high-resolution files. A visually sharp PDF cannot prove those
scientific conditions were met.
\end{scienceinpractice}

\section{Metadata}

PDF properties, portal fields, source commands, title pages, repository records,
and supplementary covers should agree on title, author names and order,
affiliations, abstract, keywords, identifiers, funders, and version.

An anonymous review PDF should not expose authors in document properties. An
identified archival PDF should use verified names. Avoid including unnecessary
personal contact details in public metadata. Inspect rather than assume what the
engine or publisher class writes.

\section{Fonts, page properties, and PDF features}

If the venue requires embedded fonts, use a PDF inspection tool to confirm that
each font is embedded or subset as allowed. Do not rely on the fact that the PDF
looks correct on the author's system. Check page size, orientation, crop boxes,
security restrictions, file size, and PDF version when specified.

Font embedding does not grant redistribution rights. A font can be technically
embedded yet contractually unsuitable. Prefer fonts distributed for the chosen
TeX workflow or otherwise licensed for the intended use.

\section{Accessibility}

Current accessibility obligations vary by venue and jurisdiction. Logical
headings, real text, meaningful link text, readable contrast, table headers, and
textual explanations of figures improve the source, but a conventional
pdf\LaTeX{} document is not automatically a fully tagged accessible PDF.

Determine the required standard and validation process. If remediation is
needed, preserve the structured source and record manual changes so an updated
document can be regenerated. Provide alt text or long descriptions through the
venue's accepted mechanism; a caption and alternative text serve related but
different purposes.

\section{Remove draft apparatus}

Search for:

\begin{terminalcode}{Typical release searches}
rg -n 'TODO|FIXME|PLACEHOLDER|DRAFT|\?\?' .
rg -n 'showkeys|linenumbers|draftwatermark' --glob '*.tex' .
\end{terminalcode}

Interpret results: a teaching document may intentionally print the word
``placeholder'', while a submission should not contain unresolved fields.
Remove line numbers, watermarks, visible changes, margin notes, draft figures,
and review comments unless the recipient explicitly requires them.

\section{Supplementary files}

For every supplement, verify title, numbering, citation from the main document,
file name, format, size, legend, and access condition. Data and code should have
a README, licence, dictionary, environment record, and execution order where
appropriate. Scan archives for hidden files, credentials, personal paths, and
restricted content.

Open the files after packaging. A valid ZIP archive can still contain the wrong
version. Ensure that supplementary equation, figure, table, and reference
numbers match the final manuscript and upload descriptions.

\section{Create a clean submission folder}

Build the folder from a release manifest. Include only required source,
bibliography, class or style files that may be redistributed, figures, and
supplements. Exclude auxiliary files, caches, old drafts, private correspondence,
credentials, raw restricted data, and unlicensed assets.

\begin{terminalcode}{Illustrative release tree}
submission_v1/
|-- main.tex
|-- sections/
|-- figures/
|-- references.bib
|-- supplementary/
|-- README.md
`-- MANIFEST.txt
\end{terminalcode}

Do not blindly flatten paths if the project uses its folder structure. Build
from the clean folder using the documented command.

\section{Inspect the final PDF}

The final PDF is the recipient's reading object even when source is supplied.
Perform automated and visual checks:

\begin{itemize}
  \item build log has no fatal, undefined-reference, or missing-file messages;
  \item page count and dimensions match expectations;
  \item fonts and images meet specified requirements;
  \item text extraction is plausible where required;
  \item every page is rendered to images and scanned for clipping or blank output;
  \item representative dense pages are inspected at full resolution;
  \item title, last page, references, appendices, and special orientations are checked.
\end{itemize}

Automation can identify candidates; a human must judge scientific and visual
meaning.

\section{Upload rehearsal}

If the submission system allows drafts, rehearse early. Portal conversion can
reorder files, replace fonts, move figures, or generate a combined PDF. Download
and inspect the system-generated proof. Compare entered metadata with the source
and secure approval from all authors when required.

Do not click the irreversible final-submission control until the responsible
author has verified the proof and declarations. Retain the confirmation and
identifier.

\begin{commonerror}
\textbf{Symptom:} the local PDF is correct, but the submission proof lacks
figures or references.

\textbf{Cause:} the portal rebuilt an incomplete source archive or interpreted
the wrong main file.

\textbf{Repair:} inspect the portal log and required file designation, correct
the clean package, rebuild locally from that exact package, and regenerate the
proof. Do not assume the portal will discover undeclared dependencies.
\end{commonerror}

\section{Archive both release and working project}

Preserve:

\begin{itemize}
  \item exact submitted PDF and source archive;
  \item supplementary files and manifest with checksums where appropriate;
  \item source-version identifier and environment record;
  \item author approvals, permission records, and submission confirmation;
  \item portal metadata or exported record;
  \item editable working project with its reproducible analysis.
\end{itemize}

The submitted release should be immutable. Continue revisions in a new version
or branch. A response to reviewers begins from what was actually submitted, not
from whichever draft happens to be open later.

\begin{goodpractice}
Use a two-person release check for high-stakes documents: one person operates the
build and upload, while another reads the checklist and verifies the evidence.
Both confirm the final proof and archive location.
\end{goodpractice}

\section{Code lab: make a clean release candidate}

\begin{codelab}
Run these commands from a copy of your project, not from your only working
copy. The result should be a rebuilt PDF, a human-readable file list, and a
checksum that identifies the exact submitted PDF.
\end{codelab}

\begin{terminalcode}{Rebuild and inspect the release candidate}
latexmk -C
latexmk -pdf main.tex
latexmk -pdf main.tex
pdffonts main.pdf
\end{terminalcode}

\begin{terminalcode}{List the files that belong in the archive}
find . -type f \
  \( -name '*.tex' -o -name '*.bib' -o -name '*.cls' \
     -o -name '*.sty' -o -name '*.bst' -o -name '*.bbx' \
     -o -name '*.cbx' -o -name '*.pdf' -o -name '*.png' \
     -o -name '*.jpg' -o -name '*.jpeg' -o -name '*.csv' \
     -o -name '*.R' -o -name 'README.md' -o -name 'latexmkrc' \) \
  | sort > MANIFEST.txt
\end{terminalcode}

\begin{terminalcode}{Record the identity of the final PDF}
shasum -a 256 main.pdf > CHECKSUMS.txt
\end{terminalcode}

Open \path{MANIFEST.txt}. Remove auxiliary products from the archive, but keep
the source, bibliography, figures, instructions, final PDF, manifest, and
checksum. On Linux, \verb|sha256sum| is the usual equivalent of the last
command.

\chapterproject

\begin{miniproject}
Use the running article to create a practice release folder and manifest. Run
its numerical check, compile from that folder, search for unresolved markers,
inspect citations and links, verify PDF properties and fonts, and complete the
included submission checklist. Record every result as pass, fail, or not
applicable with a reason; do not submit the simulated article to a journal.
\end{miniproject}

\chaptersummary

Submission is a controlled release, not merely the last compilation. A current
compliance matrix guides scientific, source, and PDF proofreading. Citations,
equations, tables, figures, metadata, fonts, accessibility, supplements, and
draft apparatus need distinct checks. The clean source package must rebuild,
the final PDF must be inspected, and the submission-system proof must be treated
as a new output. Archive the exact release and retain the editable reproducible
project for later correction or review.

\chaptercheck
\begin{chapterchecklist}
  \item A reviewed source version is frozen as the release candidate.
  \item Current venue instructions are mapped to evidence-backed checks.
  \item Scientific claims and repeated numerical results are consistent.
  \item References, equations, tables, figures, and permissions are verified.
  \item Metadata, fonts, PDF properties, links, and accessibility needs are checked.
  \item The clean folder rebuilds and contains no private or obsolete material.
  \item The portal proof and exact submitted release are inspected and archived.
\end{chapterchecklist}

\chapterexercisestitle
\begin{chapterexercises}
  \item Create a compliance matrix from ten supplied venue instructions.
  \item Design a results consistency sheet for an abstract, table, and discussion.
  \item Calculate effective resolution for three raster images at placed size.
  \item Audit a bibliography record against an authoritative source.
  \item List source-only, PDF-only, and portal-only failure modes.
  \item Build a release tree and justify each included or excluded file.
  \item Write an archive record for a fictional thesis deposit.
  \item Challenge: conduct a two-person-style submission rehearsal and document
        evidence for every checklist decision.
\end{chapterexercises}
